The way to understand this is to break down the response to the
increase in riskiness into two components. The first is the direct
precautionary effect examined in \handoutM{CRRA-RateRisk}, which works
in the way intuition suggests (more risk implies lower consumption);
the second effect is the `portfolio rebalancing' effect, examined in
\handoutA{Portfolio-CRRA}: An increase in riskiness makes the consumer
choose to invest less in the risky asset.  What \eqref{eq:PrecEffect}
tells us is that the consumer's `flight from risk' is so effective in
reducing riskiness (because the portfolio share $\riskyshare$ enters
as a {\it squared} term in \eqref{eq:PrecEffectRaw}) that the
riskiness of the {\it portfolio} actually {\it declines} as the riskiness of
the risky asset increases.  (This result was also derived in
\handoutM{Portfolio-CRRA}).  Now the overall reduction in consumption
coming through the precautionary term makes sense: Precautionary
saving is less than before {\it because the consumer optimally chooses less
  exposure to risk} than before.
